% ---------------------------------------------------------------------------
% Report skeleton. Copy to <report>/report.tex and fill in.
% Compile from the report folder:
%   pdflatex -interaction=nonstopmode -halt-on-error report.tex   (TWICE)
% Figures come from figures/make_figures.py, tables from analysis/make_tables.py.
% The section order below is the convention: DEFINITIONS BEFORE RESULTS.
% ---------------------------------------------------------------------------
\documentclass[11pt]{article}
\usepackage[margin=1in]{geometry}
\usepackage{amsmath,amssymb}
\usepackage{booktabs}
\usepackage{graphicx}
\usepackage{xcolor}
\usepackage[colorlinks=true,linkcolor=blue!60!black,urlcolor=blue!60!black]{hyperref}
\usepackage{caption}
\captionsetup{font=small,labelfont=bf}

% Macro names must be DIGIT-FREE (\arm1 is not a legal control sequence).
% \newcommand{\armA}{\textsf{baseline}}

\title{\textbf{R### --- [title]}\\
\large [subtitle: what varies, on what systems, under what protocol]\\
(\texttt{[YYYY-MM-DD]-R###-[topic]})}
\author{}
\date{[YYYY-MM-DD]}

\begin{document}
\maketitle

\begin{abstract}\noindent
[What question, what was held fixed, what varied, and the headline numbers.
State the gate status (e.g. "all gated, G1--G8 PASS") and, if a gate fails,
say so here rather than burying it. If this is a revision, say what changed
and which earlier conclusion it corrects.]
\end{abstract}

\tableofcontents

\section{What this report does, and what it does not}
% Scope, the question, and an explicit NOT-IN-SCOPE list.
% If any inherited comparison is not like-for-like with an earlier report
% (different arm, different anchor, different budget), say so HERE, loudly.

\section{Definitions --- everything used in this report}
\label{sec:defs}
% HARD REQUIREMENT. Every symbol, arm, metric, rule and constant used
% anywhere below is defined here, INCLUDING everything inherited from earlier
% work. A reader must never need another document.

\subsection{Systems, data and library}
% Equations, initial conditions, sample counts, noise model, derivative
% estimator, feature library, scaling. Say which units each stage works in.

\subsection{Methods / arms}
% One row per arm ACTUALLY USED, defined in full — not by citation.
% Composite names get each slot explained and each option enumerated.
\begin{center}\small
\begin{tabular}{llll}
\toprule
arm & [factor 1] & [factor 2] & [inner solve] \\ \midrule
\textbf{A} & & & \\
\textbf{B} & & & \\
\bottomrule
\end{tabular}
\end{center}
% If arms differ in MORE THAN ONE factor at once, either complete the
% factorial with diagnostic arms or state that no causal attribution is made.

\subsection{Protocol constants}
% Symbol -> meaning -> value -> where frozen. Separate look-alikes explicitly
% (a METHOD's cap vs a METRIC's budget; a stopping tolerance vs a threshold).
% Flag every NON-UNIFORM configuration across arms or systems and say how it
% changes the reading.

\subsection{Metrics}
% Every metric and flag, including reused ones. Name the primary metric.
% Say what is undefined (and how it is displayed) rather than plotting a 0.

\section{What was recomputed, and the staleness position}
% Table: input | FRESH or REUSED | evidence (which gate, how many rows).
% Distinguish whole-file BYTE equality from KEYED value equality.
% Note the SHA-256 ledger written by the verify script.

\section{Results}
% One subsection per axis. Every claim carries its uncertainty; paired
% differences carry paired bootstrap CIs with the resampling unit named.

\section{Gates}
% One row per gate: id, PASS/FAIL, what it asserts. Failures stay failures.

\section{Limitations}
\label{sec:lim}

\section{Recommendation}
\label{sec:reco}

\end{document}
